\chapter{Conclusions and future work}
\label{ch:conclusions}
% Purpose (THESIS_GUIDE.md, Ch. 8): answer RQ1, RQ2 and C3 directly, one
% paragraph each, then the main lessons about the method. 2-3 pages.

\section{Conclusions}
\label{sec:conclusions}
\paragraph{RQ1: comparative performance.}
\TODO{Answer from the sealed test ($H_{0,1}$, score and reward, equivalence
margin), or from the exploratory comparison with that word in the sentence.}

\paragraph{RQ2: robustness and generalizability.}
\TODO{Answer $H_{0,2}$ (degradation over amplitude) and the noise categories
(measurement and feed, by amplitude and temporal structure); failure modes and
operating envelope.}

\paragraph{C3: reliability of the synthesis.}
\TODO{Spread between the eight independent runs.}

\paragraph{The method.}
\TODO{The main lessons: the evaluation matters as much as the model; the
safeguards the loop needs; cost per run; one sentence on the tank development
work.}

\section{Future work}
\label{sec:future-work}
\begin{itemize}
\item Train the TD-MPC with feed randomisation and noise, for a symmetric
  noise comparison.
\item Train the LLM supervisor under feed fluctuations.
\item A sump-level objective that the score and the reward agree on.
\item Several operating points.
\item Compare models on the grinding circuit (Sonnet and Haiku, DeepSeek,
  effort levels).
\item Plant data, and deployment in advisory mode.
\item Retrain automatically when the regulatory loops are retuned, so that
  generated supervisors do not silently depend on the loop tuning.
\item A rule-count and readability study with operators.
\end{itemize}
\TODO{Trim and order by importance.}
